% Chapter 6 -- Throwing Things Away (Part I, Learning; AN-14 task form).
% Sources (Law 1): os/flair/finder_ops.c / finder_ops.h (a drop on the
% Trash moves the item into \TRASH, never deletes it; a name already in the
% Trash gets a numbered name; the origin is kept; a read-only item is
% refused); os/milton/kmain.c (DESK-TRASH-ICON full/empty; the Trash window
% "Trash" over \TRASH; finder_empty_trash and finder_trash_confirm: the
% Empty Trash alert, OK deletes everything in \TRASH, a locked item is not
% deleted and the Locked notice follows; the Locked notice text);
% os/flair/finder_windows.c finder_trash_alert_text (the alert wording, the
% space in K rounded up); os/flair/finder_menu.c (Empty Trash lit only while
% the Trash holds something).  Model: Macintosh User's Guide 1991 Ch. 6
% "Throwing away a document" (apple-macintosh.md Sec 8.2); System 6 guide
% pp. 83, 88 (the locked item).
% GATES: test-flair-file-ops test-flair-trash-window test-flair-empty-trash test-flair-locked-trash test-flair-locked-purge test-flair-trash-purge-views test-pubs-flair-screens
\Gates{test-flair-file-ops, test-flair-trash-window, test-flair-empty-trash,
  test-flair-locked-trash, test-flair-locked-purge,
  test-flair-trash-purge-views, test-pubs-flair-screens}
\chapter{Throwing Things Away}

To get rid of a file or a folder, you put it in the Trash and then empty
the Trash.  Until you empty it, anything in the Trash can be taken out
again.

\section{Putting something in the Trash}
\index{Trash!putting things in}\index{throwing things away}
\Operation{FINDER-TRASH}
\begin{TaskSteps}
\Step Drag the icon of the file or folder onto the Trash icon.
\Result{When the pointer is over the Trash, the Trash is darkened.}
\Step Release the button.
\Result{The icon disappears from its window, and the Trash bulges to show
that it holds something.}
\end{TaskSteps}

\FlairScreen[scale=4, label=fig:trashed]{trashed}{The Trash, holding
something}

Putting an item in the Trash does not delete it.  The item is kept in the
Trash, on the disk, and still takes up space there.  If the Trash already
holds an item of the same name, the new one is kept under a numbered name.

\section{Looking in the Trash}
\index{Trash!opening}
\Operation{FINDER-OPEN-TRASH}
\begin{TaskSteps}
\Step Double-click the Trash icon.
\Result{The Trash window opens and shows what the Trash holds.}
\end{TaskSteps}

\FlairScreen[label=fig:trashwin]{trashwin}{The Trash window}

\section{Taking something out of the Trash}
\index{Trash!taking things out}
To keep something you have put in the Trash, drag it out of the Trash
window into a Finder window, onto a folder icon, or onto the
\texttt{INITECH} icon.  It leaves the Trash and goes where you release it.

\FlairScreen[label=fig:putback]{putback}{README.TXT taken out of the Trash
and returned to Drive A Files}

\section{Emptying the Trash}
\index{Trash!emptying}\index{Empty Trash command}
\Operation{FINDER-TRASH-ALERT}
Emptying the Trash deletes everything in it from the disk, and frees the
space it used.
\begin{TaskSteps}
\Step Click a bare part of the desktop, so that the Finder is the active
application.
\Step Choose \menuitem{Empty Trash} from the \menuname{Special} menu.
\Result{An alert tells you how many items the Trash holds and how much disk
space they use, and asks you to confirm.  \menuitem{Empty Trash} is dimmed
when the Trash is empty.}
\end{TaskSteps}

\FlairScreen[label=fig:emptyask]{emptyask}{Confirming Empty Trash}

\begin{TaskSteps}[resume]
\Step Click \texttt{OK}, or press \key{Enter}.
\Result{Everything in the Trash is deleted, folders with all their
contents.  Any window showing a folder that was deleted closes.  The Trash
icon returns to normal.  To keep everything instead, click
\texttt{Cancel} or press \key{Esc}.}
\end{TaskSteps}

\begin{MacWarning}
Emptying the Trash cannot be undone.  Look in the Trash window before you
empty it.
\end{MacWarning}

\section{Locked items}
\index{locked items}
A file can be \term{locked} so that it cannot be changed or thrown away.
(InitechDOS calls a locked file \textit{read-only}; you lock and unlock
files with the \file{ATTRIB} program at the InitechDOS prompt, described in
the \doctitle{InitechDOS Reference}.)

\Operation{FINDER-LOCKED-ALERT}
If you drag a locked item to the Trash, it slides back to where it was and
an alert explains why.

\FlairScreen[label=fig:locked]{locked}{A locked file dragged to the Trash}

Click \texttt{OK}, or press \key{Enter}.  The item stays where it was.  If the Trash holds a
locked item when you empty it, FLAIR shows the same alert, and the locked
item stays in the Trash.
